\section{GPU Acceleration: Lux Accel}

\label{sec:gpu}
%% ============================================================================

\subsection{Overview}

Lux Accel is a unified GPU acceleration library that provides a single API across CUDA (NVIDIA), Metal (Apple), and WebGPU (browsers) backends. It targets cryptographic operations (NTT, polynomial multiplication, multi-scalar multiplication), AI inference kernels, and vector search distance calculations.

\subsection{Feature Matrix}

\begin{table}[H]
\centering
\caption{GPU acceleration comparison.}
\footnotesize
\begin{tabular}{@{}lccc@{}}
\toprule
Feature & \rotatebox{60}{Lux Accel} & \rotatebox{60}{Custom CUDA} & \rotatebox{60}{cuBLAS} \\
\midrule
CUDA & \yes & \yes & \yes \\
Metal & \yes & \no & \no \\
WebGPU & \yes & \no & \no \\
Unified API & \yes & \no & \pmark \\
NTT kernels & \yes & Manual & \no \\
MSM kernels & \yes & Manual & \no \\
GEMM & \yes & Manual & \yes \\
Crypto ops & \yes & Manual & \no \\
Auto-tuning & \yes & \no & \pmark \\
Rust native & \yes & \no & \no \\
Open source & \yes & N/A & \no \\
\bottomrule
\end{tabular}
\end{table}

\subsection{Performance}

NTT (Number Theoretic Transform) benchmarks, $2^{24}$ elements, Goldilocks field:

\begin{table}[H]
\centering
\caption{NTT throughput (elements per second, billions).}
\footnotesize
\begin{tabular}{@{}lrrr@{}}
\toprule
Backend & Lux Accel & Custom CUDA & icicle \\
\midrule
A100 (CUDA) & 4.2B & 3.8B & 3.5B \\
M3 Max (Metal) & 1.8B & N/A & N/A \\
RTX 4090 (CUDA) & 3.6B & 3.2B & 2.9B \\
\bottomrule
\end{tabular}
\end{table}

Lux Accel's 10--20\% advantage over hand-written CUDA comes from auto-tuned kernel launch parameters, persistent kernel execution (avoiding launch overhead for sequential operations), and memory pool pre-allocation that eliminates \texttt{cudaMalloc} calls in the critical path.

\subsection{Key Differentiator}

The unified CUDA/Metal/WebGPU API is unique. No other library provides zero-knowledge proof acceleration on Apple Silicon \textit{and} NVIDIA GPUs \textit{and} in the browser, from the same codebase. This enables a single code path for validators (NVIDIA), developers (Apple), and light clients (WebGPU).

%% ============================================================================
